\documentclass{article}
\usepackage{amsmath}
\usepackage{amsfonts}
\usepackage{amssymb}
\usepackage{geometry}
\usepackage{algorithm}
\usepackage{algorithmic}
\usepackage{graphicx}
\usepackage{tikz}
\usetikzlibrary{shapes,arrows,positioning}

\geometry{margin=1in}

\title{System2Module Forward Pass: Memory-Enhanced Reasoning Architecture}
\author{Mathematical Analysis}
\date{\today}

\begin{document}

\maketitle

\section{Notation Conventions}

\textbf{Index Notation:} We use Einstein summation convention where repeated indices are summed over, except in Hadamard products where all indices must appear on both sides to avoid implicit summation.

\textbf{Hadamard Product ($\odot$):} Element-wise multiplication between tensors of identical dimensions. In index notation, all indices appear on both sides (e.g., $A_{ij} \odot B_{ij}$).

\textbf{Broadcasting:} When dimensions don't match for Hadamard product, we use broadcasting notation (e.g., $A_{ij} \alpha_i$ means $\alpha_i$ is repeated for each column $j$).

\textbf{Softmax Representation:} For any matrix $\mathbf{M}$, we define its softmax representation as $\mathbf{M}_S = \text{softmax}(\mathbf{M}, \text{dim}=-1)$.

\textbf{Sigmoid Representation:} For any matrix $\mathbf{M}$, we define its sigmoid representation as $\mathbf{M}_\sigma = \sigma(\mathbf{M})$ where $\sigma(x) = \frac{1}{1 + e^{-x}}$.

\textbf{Function Notation:} We use $\text{RMSNorm}(\cdot)$ for RMS normalization, $\text{softmax}(\cdot)$ for softmax, and $\sigma(\cdot)$ for sigmoid consistently throughout.

\section{Overview}

The System2Module implements Kahneman's System 2 thinking through a sophisticated transformer architecture with persistent working memory. The forward pass combines input injection, multi-layer processing with memory enhancement, and reasoning history maintenance.

\section{Mathematical Formulation}

\subsection{Input Processing}

Given inputs:
\begin{align}
\mathbf{H} &\in \mathbb{R}^{B \times S \times D} \quad \text{(hidden states)} \\
\mathbf{I} &\in \mathbb{R}^{B \times S \times D} \quad \text{(input injection)} \\
\mathbf{R} &\in \mathbb{R}^{B \times P \times D} \quad \text{(reasoning history, optional)}
\end{align}

where $B$ is batch size, $S$ is sequence length, $D$ is hidden dimension, and $P$ is previous sequence length.

\subsection{Input Injection}

The first operation combines hidden states with input injection:

\begin{equation}
\mathbf{H}^{(0)} = \mathbf{H} + \mathbf{I} \in \mathbb{R}^{B \times S \times D}
\end{equation}

\textbf{Design Decision:} Simple addition is used rather than gating. This choice assumes that input injection should be additive rather than selective, allowing the model to learn appropriate combinations through subsequent layers.

\subsection{Layer-wise Processing}

For each layer $l \in \{1, 2, \ldots, L\}$ where $L$ is the number of effective layers:

\subsubsection{Working Memory Enhancement}

The working memory mechanism operates on the current state $\mathbf{x} = \mathbf{H}^{(l-1)}$:

\begin{enumerate}
\item \textbf{Memory Attention Computation:}
\begin{equation}
\mathbf{A} = \mathbf{x} \mathbf{M}^T \in \mathbb{R}^{B \times S \times M}
\end{equation}
where $\mathbf{M} \in \mathbb{R}^{M \times D}$ is the learnable memory buffer with $M$ memory slots.

\item \textbf{Attention Weight Normalization:}
\begin{equation}
\mathbf{W}_S = \text{softmax}(\mathbf{A}, \text{dim}=-1) \in \mathbb{R}^{B \times S \times M}
\end{equation}

\item \textbf{Memory Retrieval:}
\begin{equation}
\mathbf{R}_{\text{mem}} = \mathbf{W}_S \mathbf{M} \in \mathbb{R}^{B \times S \times D}
\end{equation}

\item \textbf{State Enhancement:}
\begin{equation}
\mathbf{x}_{\text{enhanced}} = \mathbf{x} + \mathbf{R}_{\text{mem}} \in \mathbb{R}^{B \times S \times D}
\end{equation}
\end{enumerate}

\textbf{Design Decision:} Memory retrieval uses attention weights rather than gating. This allows for soft, weighted combination of memory slots rather than element-wise gating. The addition operation assumes memory should be additive rather than replacing current information.

\subsubsection{Memory Retrieval Mechanisms}

The working memory supports three retrieval strategies, selected via the $\text{retrieval\_strategy}$ parameter:

\paragraph{Strategy 1: Simple Retrieval}

When $\text{retrieval\_strategy} = \text{"simple"}$:

\textbf{Parameter Definitions:}
\begin{itemize}
\item $\mathbf{M} \in \mathbb{R}^{M \times D}$: Memory matrix (learnable parameter)
\end{itemize}

\begin{enumerate}
\item \textbf{Attention Computation:}
\begin{align}
\mathbf{A}_{\text{ret}} &= \mathbf{x} \mathbf{M}^T \in \mathbb{R}^{B \times S \times M} \\
A_{\text{ret},bsi} &= \sum_{j=1}^{D} x_{bsj} M_{ij}
\end{align}

\item \textbf{Attention Weights:}
\begin{align}
\mathbf{W}_{\text{ret},S} &= \text{softmax}(\mathbf{A}_{\text{ret}}, \text{dim}=-1) \in \mathbb{R}^{B \times S \times M} \\
W_{\text{ret},S,bsi} &= \frac{\exp(A_{\text{ret},bsi})}{\sum_{i'=1}^{M} \exp(A_{\text{ret},bsi'})}
\end{align}

\item \textbf{Memory Retrieval:}
\begin{align}
\mathbf{R}_{\text{mem}} &= \mathbf{W}_{\text{ret},S} \mathbf{M} \in \mathbb{R}^{B \times S \times D} \\
R_{\text{mem},bsj} &= \sum_{i=1}^{M} W_{\text{ret},S,bsi} M_{ij}
\end{align}

\textbf{Asymptotic Complexity:} $O(B \cdot S \cdot D \cdot M)$
\begin{itemize}
\item Attention computation: $O(B \cdot S \cdot D \cdot M)$
\item Softmax computation: $O(B \cdot S \cdot M)$
\item Memory retrieval: $O(B \cdot S \cdot M \cdot D)$
\end{itemize}

\textbf{Memory Complexity:} $O(B \cdot S \cdot M + M \cdot D)$
\begin{itemize}
\item Attention matrix: $O(B \cdot S \cdot M)$
\item Memory matrix: $O(M \cdot D)$
\item Intermediate tensors: $O(B \cdot S \cdot D)$
\end{itemize}
\end{enumerate}

\paragraph{Strategy 2: Attention-Based Retrieval}

When $\text{retrieval\_strategy} = \text{"attention"}$:

\textbf{Parameter Definitions:}
\begin{itemize}
\item $\mathbf{W}_q, \mathbf{W}_k, \mathbf{W}_v \in \mathbb{R}^{D \times D}$: Query, key, value projection matrices
\item $\mathbf{W}_{\text{out}} \in \mathbb{R}^{D \times D}$: Output projection matrix
\item $\mathbf{M} \in \mathbb{R}^{M \times D}$: Memory matrix (learnable parameter)
\end{itemize}

\begin{enumerate}
\item \textbf{Multi-Head Attention:}
\begin{align}
\mathbf{R}_{\text{mem}}, \mathbf{W}_{\text{ret},S} &= \text{MultiHeadAttention}(\mathbf{x}, \mathbf{M}, \mathbf{M}) \in \mathbb{R}^{B \times S \times D} \times \mathbb{R}^{B \times S \times M} \\
\text{where: } \mathbf{Q} &= \mathbf{x} \mathbf{W}_q, \quad \mathbf{K} = \mathbf{M} \mathbf{W}_k, \quad \mathbf{V} = \mathbf{M} \mathbf{V}_v \\
R_{\text{mem},bsj} &= \sum_{h=1}^{H} \sum_{i=1}^{M} \text{softmax}\left(\frac{\sum_{k=1}^{D} Q_{bsk} K_{ik}}{\sqrt{D/H}}\right) V_{ij} \\
W_{\text{ret},S,bsi} &= \text{softmax}\left(\frac{\sum_{k=1}^{D} Q_{bsk} K_{ik}}{\sqrt{D/H}}\right)
\end{align}

\item \textbf{RMS Normalization:}
\begin{align}
\mathbf{R}_{\text{mem}} &= \text{RMSNorm}(\mathbf{R}_{\text{mem}}) \in \mathbb{R}^{B \times S \times D} \\
R_{\text{mem},bsj} &= \frac{R_{\text{mem},bsj}}{\sqrt{\frac{1}{D}\sum_{k=1}^{D} R_{\text{mem},bsk}^2 + \epsilon}} \cdot w_j
\end{align}
where $w_j$ is the learnable scale parameter and $\epsilon = 10^{-5}$ is a small constant for numerical stability.

\textbf{Asymptotic Complexity:} $O(B \cdot S \cdot D^2 + B \cdot S \cdot M \cdot D)$
\begin{itemize}
\item Attention computation: $O(B \cdot S \cdot D^2 + B \cdot S \cdot M \cdot D)$
\item RMS normalization: $O(B \cdot S \cdot D)$
\end{itemize}

\textbf{Memory Complexity:} $O(B \cdot S \cdot M + M \cdot D + B \cdot S \cdot D)$
\begin{itemize}
\item Attention weights: $O(B \cdot S \cdot M)$
\item Memory matrix: $O(M \cdot D)$
\item Query/Key/Value projections: $O(B \cdot S \cdot D)$
\item RMS normalization parameters: $O(D)$
\end{itemize}
\end{enumerate}

\paragraph{Strategy 3: Sophisticated Retrieval}

When $\text{retrieval\_strategy} = \text{"sophisticated"}$:

\textbf{Parameter Definitions:}
\begin{itemize}
\item $\boldsymbol{\alpha}_{\text{ret}} \in \mathbb{R}^{M}$: Learnable memory decay factors for retrieval
\item $\mathbf{W}_q, \mathbf{W}_k, \mathbf{W}_v \in \mathbb{R}^{D \times D}$: Query, key, value projection matrices
\item $\mathbf{W}_{\text{out}} \in \mathbb{R}^{D \times D}$: Output projection matrix
\item $\mathbf{W}_{\text{imp}} \in \mathbb{R}^{D \times 1}$: Learnable importance projection matrix
\item $\mathbf{M} \in \mathbb{R}^{M \times D}$: Memory matrix (learnable parameter)
\end{itemize}

\begin{enumerate}
\item \textbf{Memory Decay for Retrieval:}
\begin{align}
\mathbf{M}_{\text{decayed}} &= \mathbf{M} \odot (\boldsymbol{\alpha}_{\text{ret}} \mathbf{1}^T) \in \mathbb{R}^{M \times D} \\
M_{\text{decayed},ij} &= M_{ij} \alpha_{\text{ret},i} \quad \text{(broadcasting: $\alpha_{\text{ret},i}$ repeated for each column $j$)}
\end{align}

\item \textbf{Multi-Head Attention with Decayed Memory:}
\begin{align}
\mathbf{R}_{\text{mem}}, \mathbf{W}_{\text{ret},S} &= \text{MultiHeadAttention}(\mathbf{x}, \mathbf{M}_{\text{decayed}}, \mathbf{M}_{\text{decayed}}) \in \mathbb{R}^{B \times S \times D} \times \mathbb{R}^{B \times S \times M} \\
\text{where: } \mathbf{Q} &= \mathbf{x} \mathbf{W}_q, \quad \mathbf{K} = \mathbf{M}_{\text{decayed}} \mathbf{W}_k, \quad \mathbf{V} = \mathbf{M}_{\text{decayed}} \mathbf{W}_v \\
R_{\text{mem},bsj} &= \sum_{h=1}^{H} \sum_{i=1}^{M} \text{softmax}\left(\frac{\sum_{k=1}^{D} Q_{bsk} K_{ik}}{\sqrt{D/H}}\right) V_{ij} \\
W_{\text{ret},S,bsi} &= \text{softmax}\left(\frac{\sum_{k=1}^{D} Q_{bsk} K_{ik}}{\sqrt{D/H}}\right)
\end{align}

\item \textbf{RMS Normalization:}
\begin{align}
\mathbf{R}_{\text{mem}} &= \text{RMSNorm}(\mathbf{R}_{\text{mem}}) \in \mathbb{R}^{B \times S \times D} \\
R_{\text{mem},bsj} &= \frac{R_{\text{mem},bsj}}{\sqrt{\frac{1}{D}\sum_{k=1}^{D} R_{\text{mem},bsk}^2 + \epsilon}} \cdot w_j
\end{align}
where $w_j$ is the learnable scale parameter and $\epsilon = 10^{-5}$ is a small constant for numerical stability.

\item \textbf{Importance Weighting:}
\begin{align}
\mathbf{I}_{\text{weights},\sigma} &= \sigma(\mathbf{x} \mathbf{W}_{\text{imp}}) \in \mathbb{R}^{B \times S \times 1} \\
I_{\text{weights},\sigma,bs1} &= \sigma\left(\sum_{j=1}^{D} x_{bsj} W_{\text{imp},j1}\right) \\
\mathbf{R}_{\text{mem}} &= \mathbf{R}_{\text{mem}} \odot \mathbf{I}_{\text{weights},\sigma} \in \mathbb{R}^{B \times S \times D} \\
R_{\text{mem},bsj} &= R_{\text{mem},bsj} I_{\text{weights},\sigma,bs1} \quad \text{(broadcasting: $I_{\text{weights},\sigma,bs1}$ repeated for each column $j$)}
\end{align}

\textbf{Asymptotic Complexity:} $O(B \cdot S \cdot D^2 + B \cdot S \cdot M \cdot D)$
\begin{itemize}
\item Memory decay: $O(M \cdot D)$
\item Attention computation: $O(B \cdot S \cdot D^2 + B \cdot S \cdot M \cdot D)$
\item RMS normalization: $O(B \cdot S \cdot D)$
\item Importance weighting: $O(B \cdot S \cdot D)$
\end{itemize}

\textbf{Memory Complexity:} $O(B \cdot S \cdot M + M \cdot D + B \cdot S \cdot D)$
\begin{itemize}
\item Attention weights: $O(B \cdot S \cdot M)$
\item Memory matrix: $O(M \cdot D)$
\item Decayed memory: $O(M \cdot D)$
\item Query/Key/Value projections: $O(B \cdot S \cdot D)$
\item Importance weights: $O(B \cdot S)$
\item RMS normalization parameters: $O(D)$
\end{itemize}
\end{enumerate}

\textbf{Design Decisions:} 
\begin{itemize}
\item \textbf{Strategy Selection:} Three levels of sophistication allow users to choose appropriate complexity for their task.
\item \textbf{Memory Decay:} Sophisticated strategy includes separate decay factors for retrieval vs update.
\item \textbf{Importance Weighting:} Sophisticated strategy weights retrieved memory based on current state importance.
\item \textbf{Multi-Head Attention:} Attention and sophisticated strategies use multi-head attention for richer memory interactions.
\item \textbf{RMS Normalization:} Used instead of LayerNorm for better efficiency and stability, following modern LLM practices (LLaMA, PaLM, etc.).
\end{itemize}

\subsubsection{Memory Update Mechanisms}

The working memory supports three update strategies, selected via the $\text{update\_strategy}$ parameter:

\paragraph{Strategy 1: Simple Update}

When $\text{update\_strategy} = \text{"simple"}$:

\textbf{Parameter Definitions:}
\begin{itemize}
\item $\mathbf{W}_{\text{gate}} \in \mathbb{R}^{D \times M}$: Learnable gate projection matrix
\item $\mathbf{W}_{\text{update}} \in \mathbb{R}^{2D \times D}$: Learnable update projection matrix
\item $\mathbf{M} \in \mathbb{R}^{M \times D}$: Memory matrix (learnable parameter)
\end{itemize}

\begin{enumerate}
\item \textbf{Representative Selection:}
\begin{align}
\mathbf{r} &= \mathbf{x}[:, 0, :] \in \mathbb{R}^{B \times D} \\
r_{bj} &= x_{b0j} \quad \text{(first token of each sequence)}
\end{align}

\item \textbf{Update Gate Computation:}
\begin{align}
\mathbf{g}_\sigma &= \sigma(\mathbf{r} \mathbf{W}_{\text{gate}}) \in \mathbb{R}^{B \times M} \\
g_{\sigma,bi} &= \sigma\left(\sum_{j=1}^{D} r_{bj} W_{\text{gate},ji}\right) \quad \text{where } \sigma(x) = \frac{1}{1 + e^{-x}}
\end{align}

\item \textbf{Memory Update Preparation:}
\begin{align}
\bar{\mathbf{m}} &= \frac{1}{M}\sum_{i=1}^{M} \mathbf{M}[i, :] \in \mathbb{R}^{D} \\
\bar{m}_j &= \frac{1}{M}\sum_{i=1}^{M} M_{ij} \\
\mathbf{u} &= \mathbf{W}_{\text{update}} \begin{bmatrix} \mathbf{r} \\ \bar{\mathbf{m}} \end{bmatrix} \in \mathbb{R}^{B \times D} \\
u_{bj} &= \sum_{k=1}^{2D} W_{\text{update},kj} \begin{bmatrix} r_{bk} \\ \bar{m}_k \end{bmatrix}
\end{align}

\item \textbf{Vectorized Memory Update:}
\begin{align}
\bar{\mathbf{g}}_\sigma &= \frac{1}{B}\sum_{b=1}^{B} \mathbf{g}_\sigma[b, :] \in \mathbb{R}^{M} \\
\bar{g}_{\sigma,i} &= \frac{1}{B}\sum_{b=1}^{B} g_{\sigma,bi} \\
\bar{\mathbf{u}} &= \frac{1}{B}\sum_{b=1}^{B} \mathbf{u}[b, :] \in \mathbb{R}^{D} \\
\bar{u}_j &= \frac{1}{B}\sum_{b=1}^{B} u_{bj} \\
\mathbf{M} &\leftarrow (1 - \bar{\mathbf{g}}_\sigma \mathbf{1}^T) \odot \mathbf{M} + (\bar{\mathbf{g}}_\sigma \mathbf{1}^T) \odot (\mathbf{1} \bar{\mathbf{u}}^T) \\
M_{ij} &\leftarrow (1 - \bar{g}_{\sigma,i}) M_{ij} + \bar{g}_{\sigma,i} \bar{u}_j \quad \text{(broadcasting: $\bar{g}_{\sigma,i}$ and $\bar{u}_j$ create $M \times D$ matrices)}
\end{align}

\textbf{Asymptotic Complexity:} $O(B \cdot D \cdot M + B \cdot D^2)$
\begin{itemize}
\item Gate computation: $O(B \cdot D \cdot M)$
\item Update preparation: $O(B \cdot D^2 + M \cdot D)$
\item Memory update: $O(M \cdot D)$
\end{itemize}

\textbf{Memory Complexity:} $O(B \cdot M + M \cdot D + B \cdot D)$
\begin{itemize}
\item Gate matrix: $O(B \cdot M)$
\item Memory matrix: $O(M \cdot D)$
\item Update signal: $O(B \cdot D)$
\item Gate projection matrix: $O(D \cdot M)$
\item Update projection matrix: $O(2D \cdot D)$
\end{itemize}
\end{enumerate}

\paragraph{Strategy 2: Attention-Based Update}

When $\text{update\_strategy} = \text{"attention"}$:

\textbf{Parameter Definitions:}
\begin{itemize}
\item $\mathbf{W}_q, \mathbf{W}_k, \mathbf{W}_v \in \mathbb{R}^{D \times D}$: Query, key, value projection matrices
\item $\mathbf{W}_{\text{out}} \in \mathbb{R}^{D \times D}$: Output projection matrix
\item $\mathbf{M} \in \mathbb{R}^{M \times D}$: Memory matrix (learnable parameter)
\end{itemize}

\begin{enumerate}
\item \textbf{Memory Attention:}
\begin{align}
\mathbf{A}_{\text{mem}}, \mathbf{W}_{\text{attn},S} &= \text{MultiHeadAttention}(\mathbf{x}, \mathbf{M}, \mathbf{M}) \in \mathbb{R}^{B \times S \times D} \times \mathbb{R}^{B \times S \times M} \\
\text{where: } \mathbf{Q} &= \mathbf{x} \mathbf{W}_q, \quad \mathbf{K} = \mathbf{M} \mathbf{W}_k, \quad \mathbf{V} = \mathbf{M} \mathbf{W}_v \\
A_{\text{mem},bsj} &= \sum_{h=1}^{H} \sum_{i=1}^{M} \text{softmax}\left(\frac{\sum_{k=1}^{D} Q_{bsk} K_{ik}}{\sqrt{D/H}}\right) V_{ij} \\
W_{\text{attn},S,bsi} &= \text{softmax}\left(\frac{\sum_{k=1}^{D} Q_{bsk} K_{ik}}{\sqrt{D/H}}\right)
\end{align}

\item \textbf{Update Signal:}
\begin{align}
\mathbf{s} &= \frac{1}{S}\sum_{s=1}^{S} \mathbf{A}_{\text{mem}}[:, s, :] \in \mathbb{R}^{B \times D} \\
s_{bj} &= \frac{1}{S}\sum_{s=1}^{S} A_{\text{mem},bsj}
\end{align}

\item \textbf{Attention-Based Gates:}
\begin{align}
\mathbf{g}_{\text{attn},\sigma} &= \sigma\left(\frac{1}{S}\sum_{s=1}^{S} \mathbf{W}_{\text{attn},S}[:, s, :]\right) \in \mathbb{R}^{B \times M} \\
g_{\text{attn},\sigma,bi} &= \sigma\left(\frac{1}{S}\sum_{s=1}^{S} W_{\text{attn},S,bsi}\right)
\end{align}

\item \textbf{Vectorized Memory Update:}
\begin{align}
\bar{\mathbf{g}}_{\text{attn},\sigma} &= \frac{1}{B}\sum_{b=1}^{B} \mathbf{g}_{\text{attn},\sigma}[b, :] \in \mathbb{R}^{M} \\
\bar{g}_{\text{attn},\sigma,i} &= \frac{1}{B}\sum_{b=1}^{B} g_{\text{attn},\sigma,bi} \\
\bar{\mathbf{s}} &= \frac{1}{B}\sum_{b=1}^{B} \mathbf{s}[b, :] \in \mathbb{R}^{D} \\
\bar{s}_j &= \frac{1}{B}\sum_{b=1}^{B} s_{bj} \\
\mathbf{M} &\leftarrow (1 - \bar{\mathbf{g}}_{\text{attn},\sigma} \mathbf{1}^T) \odot \mathbf{M} + (\bar{\mathbf{g}}_{\text{attn},\sigma} \mathbf{1}^T) \odot (\mathbf{1} \bar{\mathbf{s}}^T) \\
M_{ij} &\leftarrow (1 - \bar{g}_{\text{attn},\sigma,i}) M_{ij} + \bar{g}_{\text{attn},\sigma,i} \bar{s}_j \quad \text{(broadcasting: $\bar{g}_{\text{attn},\sigma,i}$ and $\bar{s}_j$ create $M \times D$ matrices)}
\end{align}

\textbf{Asymptotic Complexity:} $O(B \cdot S \cdot D^2 + B \cdot S \cdot M \cdot D)$
\begin{itemize}
\item Attention computation: $O(B \cdot S \cdot D^2 + B \cdot S \cdot M \cdot D)$
\item Update signal: $O(B \cdot S \cdot D)$
\item Memory update: $O(M \cdot D)$
\end{itemize}

\textbf{Memory Complexity:} $O(B \cdot S \cdot M + M \cdot D + B \cdot S \cdot D)$
\begin{itemize}
\item Attention weights: $O(B \cdot S \cdot M)$
\item Memory matrix: $O(M \cdot D)$
\item Attention output: $O(B \cdot S \cdot D)$
\item Query/Key/Value projections: $O(B \cdot S \cdot D)$
\item Gate matrix: $O(B \cdot M)$
\end{itemize}
\end{enumerate}

\paragraph{Strategy 3: Sophisticated Update}

When $\text{update\_strategy} = \text{"sophisticated"}$:

\textbf{Parameter Definitions:}
\begin{itemize}
\item $\boldsymbol{\alpha} \in \mathbb{R}^{M}$: Learnable memory decay factors
\item $\mathbf{W}_{\text{imp}} \in \mathbb{R}^{D \times 1}$: Learnable importance projection matrix
\item $\mathbf{W}_q, \mathbf{W}_k, \mathbf{W}_v \in \mathbb{R}^{D \times D}$: Query, key, value projection matrices
\item $\mathbf{W}_{\text{out}} \in \mathbb{R}^{D \times D}$: Output projection matrix
\item $\mathbf{W}_{\text{gate}} \in \mathbb{R}^{D \times M}$: Learnable gate projection matrix
\item $\mathbf{M} \in \mathbb{R}^{M \times D}$: Memory matrix (learnable parameter)
\end{itemize}

\begin{enumerate}
\item \textbf{Memory Decay:}
\begin{align}
\mathbf{M} &\leftarrow \mathbf{M} \odot (\boldsymbol{\alpha} \mathbf{1}^T) \\
M_{ij} &\leftarrow M_{ij} \alpha_i \quad \text{(broadcasting: $\alpha_i$ repeated for each column $j$)}
\end{align}
where $\mathbf{1} \in \mathbb{R}^{D}$ is a vector of ones, making $\boldsymbol{\alpha} \mathbf{1}^T \in \mathbb{R}^{M \times D}$ for proper Hadamard product dimensions.

\item \textbf{Importance-Weighted Representative:}
\begin{align}
\mathbf{w}_S &= \text{softmax}(\sigma(\mathbf{x} \mathbf{W}_{\text{imp}}), \text{dim}=1) \in \mathbb{R}^{B \times S} \\
w_{S,bs} &= \frac{\exp(\sigma(\sum_{j=1}^{D} x_{bsj} W_{\text{imp},j1}))}{\sum_{s'=1}^{S} \exp(\sigma(\sum_{j=1}^{D} x_{bs'j} W_{\text{imp},j1}))} \\
\mathbf{r}_{\text{imp}} &= \sum_{s=1}^{S} \mathbf{w}_S[:, s] \odot \mathbf{x}[:, s, :] \in \mathbb{R}^{B \times D} \\
r_{\text{imp},bj} &= \sum_{s=1}^{S} w_{S,bs} x_{bsj} \quad \text{(broadcasting: $w_{S,bs}$ repeated for each column $j$)}
\end{align}

\textbf{Role of $\mathbf{w}_S$ and $\mathbf{r}_{\text{imp}}$:}
\begin{itemize}
\item \textbf{$\mathbf{w}_S$ (importance weights):} Softmax-normalized attention weights that determine how much each sequence position contributes to the representative. Each row sums to 1, creating a probability distribution over sequence positions.
\item \textbf{$\mathbf{r}_{\text{imp}}$ (importance-weighted representative):} A single representative vector per batch item that combines all sequence positions according to their learned importance. This replaces the simple "first token" approach with a learned, context-aware selection mechanism.
\end{itemize}

\item \textbf{Attention-Based Memory Interaction:}
\begin{align}
\mathbf{A}_{\text{soph}}, \mathbf{W}_{\text{soph},S} &= \text{MultiHeadAttention}(\mathbf{r}_{\text{imp}}, \mathbf{M}, \mathbf{M}) \in \mathbb{R}^{B \times 1 \times D} \times \mathbb{R}^{B \times 1 \times M} \\
\text{where: } \mathbf{Q} &= \mathbf{r}_{\text{imp}} \mathbf{W}_q, \quad \mathbf{K} = \mathbf{M} \mathbf{W}_k, \quad \mathbf{V} = \mathbf{M} \mathbf{W}_v \\
A_{\text{soph},b1j} &= \sum_{h=1}^{H} \sum_{i=1}^{M} \text{softmax}\left(\frac{\sum_{k=1}^{D} Q_{b1k} K_{ik}}{\sqrt{D/H}}\right) V_{ij} \\
W_{\text{soph},S,b1i} &= \text{softmax}\left(\frac{\sum_{k=1}^{D} Q_{b1k} K_{ik}}{\sqrt{D/H}}\right)
\end{align}

\item \textbf{Adaptive Learning Rates:}
\begin{align}
\mathbf{g}_{\text{base},\sigma} &= \sigma(\mathbf{r}_{\text{imp}} \mathbf{W}_{\text{gate}}) \in \mathbb{R}^{B \times M} \\
g_{\text{base},\sigma,bi} &= \sigma\left(\sum_{j=1}^{D} r_{\text{imp},bj} W_{\text{gate},ji}\right) \\
\mathbf{g}_{\text{final}} &= \mathbf{g}_{\text{base},\sigma} \odot \mathbf{W}_{\text{soph},S}[:, 0, :] \in \mathbb{R}^{B \times M} \\
g_{\text{final},bi} &= g_{\text{base},\sigma,bi} W_{\text{soph},S,b0i} \quad \text{(true Hadamard product: all indices $b,i$ on both sides)} \\
\boldsymbol{\lambda} &= \mathbf{g}_{\text{final}} \odot (\mathbf{1} (1 - \boldsymbol{\alpha})^T) \in \mathbb{R}^{B \times M} \\
\lambda_{bi} &= g_{\text{final},bi} (1 - \alpha_i) \quad \text{(broadcasting: $\alpha_i$ repeated for each batch $b$)}
\end{align}

\item \textbf{Vectorized Sophisticated Update:}
\begin{align}
\bar{\boldsymbol{\lambda}} &= \frac{1}{B}\sum_{b=1}^{B} \boldsymbol{\lambda}[b, :] \in \mathbb{R}^{M} \\
\bar{\lambda}_i &= \frac{1}{B}\sum_{b=1}^{B} \lambda_{bi} \\
\bar{\mathbf{u}}_{\text{soph}} &= \frac{1}{B}\sum_{b=1}^{B} \mathbf{u}[b, :] \in \mathbb{R}^{D} \\
\bar{u}_{\text{soph},j} &= \frac{1}{B}\sum_{b=1}^{B} u_{bj} \\
\mathbf{M} &\leftarrow (1 - \bar{\boldsymbol{\lambda}} \mathbf{1}^T) \odot \mathbf{M} + (\bar{\boldsymbol{\lambda}} \mathbf{1}^T) \odot (\mathbf{1} \bar{\mathbf{u}}_{\text{soph}}^T) \\
M_{ij} &\leftarrow (1 - \bar{\lambda}_i) M_{ij} + \bar{\lambda}_i \bar{u}_{\text{soph},j} \quad \text{(broadcasting: $\bar{\lambda}_i$ and $\bar{u}_{\text{soph},j}$ create $M \times D$ matrices)}
\end{align}

\textbf{Asymptotic Complexity:} $O(B \cdot S \cdot D + B \cdot D^2 + B \cdot M \cdot D)$
\begin{itemize}
\item Importance computation: $O(B \cdot S \cdot D)$
\item Attention computation: $O(B \cdot D^2 + B \cdot M \cdot D)$
\item Gate computation: $O(B \cdot D \cdot M)$
\item Memory update: $O(M \cdot D)$
\end{itemize}

\textbf{Memory Complexity:} $O(B \cdot S + M \cdot D + B \cdot D + B \cdot M)$
\begin{itemize}
\item Importance weights: $O(B \cdot S)$
\item Memory matrix: $O(M \cdot D)$
\item Attention weights: $O(B \cdot M)$
\item Query/Key/Value projections: $O(B \cdot D)$
\item Gate matrices: $O(B \cdot M)$
\item Decay factors: $O(M)$
\end{itemize}
\end{enumerate}

\textbf{Design Decisions:} 
\begin{itemize}
\item \textbf{Strategy Selection:} Three levels of sophistication allow users to choose appropriate complexity for their task.
\item \textbf{Vectorization:} All strategies use batch-averaged updates for computational efficiency.
\item \textbf{Representative Selection:} Simple uses first token, sophisticated uses importance-weighted average.
\item \textbf{Memory Decay:} Only sophisticated strategy includes forgetting mechanisms.
\end{itemize}

\subsubsection{Sophisticated Attention}

The attention mechanism processes the memory-enhanced state:

\begin{enumerate}
\item \textbf{Query, Key, Value Projections:}
\begin{align}
\mathbf{Q} &= \mathbf{W}_q \mathbf{x}_{\text{enhanced}} \in \mathbb{R}^{B \times S \times D} \\
\mathbf{K} &= \mathbf{W}_k \mathbf{x}_{\text{enhanced}} \in \mathbb{R}^{B \times S \times D} \\
\mathbf{V} &= \mathbf{W}_v \mathbf{x}_{\text{enhanced}} \in \mathbb{R}^{B \times S \times D}
\end{align}

\item \textbf{Multi-Head Reshaping:}
\begin{align}
\mathbf{Q} &\in \mathbb{R}^{B \times H \times S \times D_h} \\
\mathbf{K} &\in \mathbb{R}^{B \times H \times S \times D_h} \\
\mathbf{V} &\in \mathbb{R}^{B \times H \times S \times D_h}
\end{align}
where $H$ is the number of heads and $D_h = D/H$.

\item \textbf{Rotary Position Embedding:}
\begin{align}
\mathbf{Q}' &= \mathbf{Q} \odot \cos(\boldsymbol{\Theta}) + \text{rotate}(\mathbf{Q}) \odot \sin(\boldsymbol{\Theta}) \\
\mathbf{K}' &= \mathbf{K} \odot \cos(\boldsymbol{\Theta}) + \text{rotate}(\mathbf{K}) \odot \sin(\boldsymbol{\Theta})
\end{align}
where $\boldsymbol{\Theta} \in \mathbb{R}^{S \times D_h}$ contains position-dependent frequencies.

\item \textbf{Scaled Dot-Product Attention:}
\begin{equation}
\mathbf{A}_{\text{attn},S} = \text{softmax}\left(\frac{\mathbf{Q}' \mathbf{K}'^T}{\sqrt{D_h}}\right) \in \mathbb{R}^{B \times H \times S \times S}
\end{equation}

\item \textbf{Attention Output:}
\begin{equation}
\mathbf{O}_{\text{attn}} = \mathbf{A}_{\text{attn},S} \mathbf{V} \in \mathbb{R}^{B \times H \times S \times D_h}
\end{equation}

\item \textbf{Reshape and Project:}
\begin{equation}
\mathbf{O}_{\text{attn}} \in \mathbb{R}^{B \times S \times D} \rightarrow \mathbf{W}_{\text{out}} \mathbf{O}_{\text{attn}} \in \mathbb{R}^{B \times S \times D}
\end{equation}

\item \textbf{Reasoning History Attention (if available):}
\begin{equation}
\mathbf{O}_{\text{hist}}, \_ = \text{MultiHeadAttention}(\mathbf{O}_{\text{attn}}, \mathbf{R}, \mathbf{R}) \in \mathbb{R}^{B \times S \times D}
\end{equation}

\item \textbf{Final Attention Output:}
\begin{equation}
\mathbf{O}_{\text{final}} = \mathbf{O}_{\text{attn}} + \mathbf{O}_{\text{hist}} \in \mathbb{R}^{B \times S \times D}
\end{equation}
\end{enumerate}

\textbf{Design Decision:} Reasoning history is combined via addition rather than concatenation or gating. This maintains dimensional consistency while allowing the model to learn appropriate combinations through residual connections.

\subsubsection{Residual Connection and Normalization}

\begin{equation}
\mathbf{x}_{\text{attn}} = \mathbf{x}_{\text{enhanced}} + \mathbf{O}_{\text{final}} \in \mathbb{R}^{B \times S \times D}
\end{equation}

\begin{equation}
\mathbf{x}_{\text{norm1}} = \text{RMSNorm}(\mathbf{x}_{\text{attn}}) \in \mathbb{R}^{B \times S \times D}
\end{equation}

where RMSNorm is defined as:
\begin{equation}
\text{RMSNorm}(\mathbf{x}) = \mathbf{x} \odot \frac{1}{\sqrt{\frac{1}{D}\sum_{i=1}^{D} x_i^2 + \epsilon}}
\end{equation}

\subsubsection{MLP Processing (SwiGLU)}

\begin{enumerate}
\item \textbf{Gate and Up Projections:}
\begin{align}
\mathbf{G} &= \mathbf{W}_{\text{gate}} \mathbf{x}_{\text{norm1}} \in \mathbb{R}^{B \times S \times D_{\text{int}}} \\
\mathbf{U} &= \mathbf{W}_{\text{up}} \mathbf{x}_{\text{norm1}} \in \mathbb{R}^{B \times S \times D_{\text{int}}}
\end{align}

\item \textbf{Gated Activation:}
\begin{equation}
\mathbf{O}_{\text{mlp}} = \mathbf{W}_{\text{down}} (\text{SiLU}(\mathbf{G}) \odot \mathbf{U}) \in \mathbb{R}^{B \times S \times D}
\end{equation}

\item \textbf{Residual Connection and Normalization:}
\begin{align}
\mathbf{x}_{\text{mlp}} &= \mathbf{x}_{\text{norm1}} + \mathbf{O}_{\text{mlp}} \in \mathbb{R}^{B \times S \times D} \\
\mathbf{x}_{\text{final}} &= \text{RMSNorm}(\mathbf{x}_{\text{mlp}}) \in \mathbb{R}^{B \times S \times D}
\end{align}
\end{enumerate}

\textbf{Design Decision:} SwiGLU uses element-wise multiplication ($\odot$) rather than addition for gating. This allows for more sophisticated non-linear combinations and has shown superior performance in modern language models.

\subsection{Reasoning History Update}

After processing all layers, if in training mode:

\begin{equation}
\mathbf{R}_{\text{new}} = \begin{cases}
\mathbf{H}^{(L)} & \text{if } \mathbf{R} = \text{None} \\
\text{concat}(\mathbf{R}, \mathbf{H}^{(L)})[:, -P_{\max}:, :] & \text{otherwise}
\end{cases}
\end{equation}

where $P_{\max}$ is the maximum history length.

\textbf{Design Decision:} Concatenation is used for history accumulation rather than averaging or gating. This preserves the temporal sequence of reasoning steps, allowing the model to attend to specific historical positions.

\section{Algorithm Summary}

\begin{algorithm}
\caption{System2Module Forward Pass with Memory Update and Retrieval Strategies}
\begin{algorithmic}[1]
\STATE \textbf{Input:} $\mathbf{H} \in \mathbb{R}^{B \times S \times D}$, $\mathbf{I} \in \mathbb{R}^{B \times S \times D}$, $\mathbf{R} \in \mathbb{R}^{B \times P \times D}$, $\text{update\_strategy} \in \{\text{"simple"}, \text{"attention"}, \text{"sophisticated"}\}$, $\text{retrieval\_strategy} \in \{\text{"simple"}, \text{"attention"}, \text{"sophisticated"}\}$
\STATE $\mathbf{H}^{(0)} \leftarrow \mathbf{H} + \mathbf{I}$
\FOR{$l = 1$ to $L$}
    \STATE $\mathbf{x} \leftarrow \mathbf{H}^{(l-1)}$
    \STATE $\mathbf{x}_{\text{enhanced}} \leftarrow \text{WorkingMemory}(\mathbf{x}, \text{update\_strategy}, \text{retrieval\_strategy})$
    \STATE $\mathbf{x}_{\text{attn}} \leftarrow \mathbf{x}_{\text{enhanced}} + \text{Attention}(\mathbf{x}_{\text{enhanced}}, \mathbf{R})$
    \STATE $\mathbf{x}_{\text{norm1}} \leftarrow \text{RMSNorm}(\mathbf{x}_{\text{attn}})$
    \STATE $\mathbf{x}_{\text{mlp}} \leftarrow \mathbf{x}_{\text{norm1}} + \text{SwiGLU}(\mathbf{x}_{\text{norm1}})$
    \STATE $\mathbf{H}^{(l)} \leftarrow \text{RMSNorm}(\mathbf{x}_{\text{mlp}})$
\ENDFOR
\IF{training}
    \STATE $\mathbf{R} \leftarrow \text{UpdateHistory}(\mathbf{H}^{(L)})$
\ENDIF
\STATE \textbf{Return:} $\mathbf{H}^{(L)}$
\end{algorithmic}
\end{algorithm}

\begin{algorithm}
\caption{WorkingMemory with Vectorized Updates and Retrieval}
\begin{algorithmic}[1]
\STATE \textbf{Input:} $\mathbf{x} \in \mathbb{R}^{B \times S \times D}$, $\text{update\_strategy}$, $\text{retrieval\_strategy}$, $\text{update\_memory} = \text{True}$
\IF{$\text{retrieval\_strategy} = \text{"simple"}$}
    \STATE $\mathbf{A} \leftarrow \mathbf{x} \mathbf{M}^T$ \COMMENT{Memory attention}
    \STATE $\mathbf{W}_S \leftarrow \text{softmax}(\mathbf{A}, \text{dim}=-1)$ \COMMENT{Attention weights}
    \STATE $\mathbf{R}_{\text{mem}} \leftarrow \mathbf{W}_S \mathbf{M}$ \COMMENT{Memory retrieval}
\ELSIF{$\text{retrieval\_strategy} = \text{"attention"}$}
    \STATE $\mathbf{R}_{\text{mem}}, \mathbf{W}_S \leftarrow \text{MultiHeadAttention}(\mathbf{x}, \mathbf{M}, \mathbf{M})$ \COMMENT{Multi-head attention retrieval}
    \STATE $\mathbf{R}_{\text{mem}} \leftarrow \text{RMSNorm}(\mathbf{R}_{\text{mem}})$ \COMMENT{RMS normalization}
\ELSIF{$\text{retrieval\_strategy} = \text{"sophisticated"}$}
    \STATE $\mathbf{M}_{\text{decayed}} \leftarrow \mathbf{M} \odot (\boldsymbol{\alpha}_{\text{ret}} \mathbf{1}^T)$ \COMMENT{Memory decay for retrieval}
    \STATE $\mathbf{R}_{\text{mem}}, \mathbf{W}_S \leftarrow \text{MultiHeadAttention}(\mathbf{x}, \mathbf{M}_{\text{decayed}}, \mathbf{M}_{\text{decayed}})$ \COMMENT{Multi-head attention with decayed memory}
    \STATE $\mathbf{R}_{\text{mem}} \leftarrow \text{RMSNorm}(\mathbf{R}_{\text{mem}})$ \COMMENT{RMS normalization}
    \STATE $\mathbf{I}_{\text{weights},\sigma} \leftarrow \sigma(\mathbf{x} \mathbf{W}_{\text{imp}})$ \COMMENT{Importance weights}
    \STATE $\mathbf{R}_{\text{mem}} \leftarrow \mathbf{R}_{\text{mem}} \odot \mathbf{I}_{\text{weights},\sigma}$ \COMMENT{Importance weighting}
\ENDIF
\STATE $\mathbf{x}_{\text{enhanced}} \leftarrow \mathbf{x} + \mathbf{R}_{\text{mem}}$ \COMMENT{State enhancement}
\IF{update\_memory}
    \IF{update\_strategy = "simple"}
        \STATE $\mathbf{r} \leftarrow \mathbf{x}[:, 0, :]$ \COMMENT{First token representative}
        \STATE $\mathbf{g}_\sigma \leftarrow \sigma(\mathbf{W}_{\text{gate}} \mathbf{r})$ \COMMENT{Update gates}
        \STATE $\mathbf{u} \leftarrow \mathbf{W}_{\text{update}} [\mathbf{r}; \bar{\mathbf{m}}]$ \COMMENT{Update signal}
        \STATE $\mathbf{M} \leftarrow (1 - \bar{\mathbf{g}}_\sigma \mathbf{1}^T) \odot \mathbf{M} + (\bar{\mathbf{g}}_\sigma \mathbf{1}^T) \odot (\mathbf{1} \bar{\mathbf{u}}^T)$ \COMMENT{Vectorized update}
    \ELSIF{update\_strategy = "attention"}
        \STATE $(\mathbf{A}_{\text{mem}}, \mathbf{W}_{\text{attn},S}) \leftarrow \text{MultiHeadAttention}(\mathbf{x}, \mathbf{M}, \mathbf{M})$
        \STATE $\mathbf{s} \leftarrow \frac{1}{S}\sum_{s=1}^{S} \mathbf{A}_{\text{mem}}[:, s, :]$ \COMMENT{Update signal}
        \STATE $\mathbf{g}_{\text{attn},\sigma} \leftarrow \sigma(\frac{1}{S}\sum_{s=1}^{S} \mathbf{W}_{\text{attn},S}[:, s, :])$ \COMMENT{Attention gates}
        \STATE $\mathbf{M} \leftarrow (1 - \bar{\mathbf{g}}_{\text{attn},\sigma} \mathbf{1}^T) \odot \mathbf{M} + (\bar{\mathbf{g}}_{\text{attn},\sigma} \mathbf{1}^T) \odot (\mathbf{1} \bar{\mathbf{s}}^T)$ \COMMENT{Vectorized update}
    \ELSIF{update\_strategy = "sophisticated"}
        \STATE $\mathbf{M} \leftarrow \mathbf{M} \odot (\boldsymbol{\alpha} \mathbf{1}^T)$ \COMMENT{Memory decay (row-wise scaling)}
        \STATE $\mathbf{w}_S \leftarrow \text{softmax}(\sigma(\mathbf{W}_{\text{imp}} \mathbf{x}), \text{dim}=1)$ \COMMENT{Importance weights}
        \STATE $\mathbf{r}_{\text{imp}} \leftarrow \sum_{s=1}^{S} \mathbf{w}_S[:, s] \odot \mathbf{x}[:, s, :]$ \COMMENT{Importance representative (broadcasting: $w_{S,bs}$ repeated for each column)}
        \STATE $(\mathbf{A}_{\text{soph}}, \mathbf{W}_{\text{soph},S}) \leftarrow \text{MultiHeadAttention}(\mathbf{r}_{\text{imp}}, \mathbf{M}, \mathbf{M})$
        \STATE $\boldsymbol{\lambda} \leftarrow \sigma(\mathbf{W}_{\text{gate}} \mathbf{r}_{\text{imp}}) \odot \mathbf{W}_{\text{soph},S}[:, 0, :] \odot (\mathbf{1} (1 - \boldsymbol{\alpha})^T)$ \COMMENT{Adaptive learning rates}
        \STATE $\mathbf{M} \leftarrow (1 - \bar{\boldsymbol{\lambda}} \mathbf{1}^T) \odot \mathbf{M} + (\bar{\boldsymbol{\lambda}} \mathbf{1}^T) \odot (\mathbf{1} \bar{\mathbf{u}}^T)$ \COMMENT{Vectorized sophisticated update}
    \ENDIF
\ENDIF
\STATE \textbf{Return:} $\mathbf{x}_{\text{enhanced}}$
\end{algorithmic}
\end{algorithm}

\section{Design Analysis}

\subsection{Combination Strategies}

\begin{itemize}
\item \textbf{Addition vs. Concatenation:} Addition is used for residual connections and memory enhancement, maintaining dimensional consistency while allowing learned combinations. Concatenation is used for history accumulation to preserve temporal information.

\item \textbf{Gating vs. Attention:} Memory retrieval uses attention for soft, weighted combinations. MLP processing uses gating (SwiGLU) for non-linear feature interactions.

\item \textbf{Representative Selection:} First token selection prioritizes efficiency over information content, relying on the loss function to adapt the choice.

\end{itemize}

\subsection{Memory Architecture}

The working memory mechanism implements a dynamic memory system with three update strategies:

\subsubsection{Strategy Comparison}

\begin{table}[h]
\centering
\begin{tabular}{|l|c|c|c|}
\hline
\textbf{Feature} & \textbf{Simple} & \textbf{Attention} & \textbf{Sophisticated} \\
\hline
Representative Selection & First token & Attention-weighted & Importance-weighted \\
Memory Decay & No & No & Yes \\
Attention Mechanism & No & Yes & Yes \\
Adaptive Learning Rates & No & No & Yes \\
Computational Cost & Low & Medium & High \\
Memory Change Magnitude & High & Medium & Low \\
\hline
\end{tabular}
\caption{Comparison of memory update strategies}
\end{table}

\subsubsection{Vectorization Benefits}

All strategies use vectorized operations for computational efficiency:
\begin{itemize}
\item \textbf{Batch Averaging:} Update signals are averaged across batch dimension
\item \textbf{Element-wise Operations:} Memory updates use broadcasting for efficiency
\item \textbf{No Nested Loops:} Eliminated inefficient Python loops
\item \textbf{GPU Acceleration:} Operations can run in parallel on GPU
\end{itemize}

\subsubsection{Retrieval Strategy Comparison}

\begin{table}[h]
\centering
\begin{tabular}{|l|c|c|c|}
\hline
\textbf{Feature} & \textbf{Simple} & \textbf{Attention} & \textbf{Sophisticated} \\
\hline
Attention Mechanism & Basic & Multi-head & Multi-head \\
Memory Decay & No & No & Yes \\
Importance Weighting & No & No & Yes \\
RMS Normalization & No & Yes & Yes \\
Computational Cost & Low & Medium & High \\
Retrieval Quality & Basic & Rich & Adaptive \\
\hline
\end{tabular}
\caption{Comparison of memory retrieval strategies}
\end{table}

\subsubsection{Performance Characteristics}

\begin{itemize}
\item \textbf{Simple Strategy:} Fastest execution, largest memory changes, minimal parameters
\item \textbf{Attention Strategy:} Balanced performance, uses entire sequence information
\item \textbf{Sophisticated Strategy:} Most controlled updates, includes forgetting mechanisms
\end{itemize}

\subsubsection{Complexity Comparison}

\begin{table}[h]
\centering
\small
\hyphenpenalty=10000
\begin{tabular}{|p{2.2cm}|p{3.5cm}|p{3.5cm}|p{2.8cm}|}
\hline
\textbf{Strategy} & \textbf{Time Complexity} & \textbf{Memory Complexity} & \textbf{Dominant Factor} \\
\hline
Simple Retrieval & $O(BSDM)$ & $O(BSM + MD)$ & Memory slots $\times$ dims \\
\hline
Attention Retrieval & $O(BSD^2 + BSMD)$ & $O(BSM + MD + BSD)$ & Attention computation \\
\hline
Sophisticated Retrieval & $O(BSD^2 + BSMD)$ & $O(BSM + MD + BSD)$ & Attention + importance \\
\hline
Simple Update & $O(BDM + BD^2)$ & $O(BM + MD + BD)$ & Gate computation \\
\hline
Attention Update & $O(BSD^2 + BSMD)$ & $O(BSM + MD + BSD)$ & Attention computation \\
\hline
Sophisticated Update & $O(BSD + BD^2 + BMD)$ & $O(BS + MD + BD + BM)$ & Importance + attention \\
\hline
\end{tabular}
\caption{Computational and memory complexity comparison across all strategies}
\end{table}

This design provides a flexible foundation for persistent memory with configurable sophistication levels while maintaining computational efficiency through vectorization.

\section{Vectorization Improvements}

\subsection{Problem with Original Implementation}

The original memory update implementation used nested loops:
\begin{verbatim}
for b in range(batch_size):
    for m in range(memory_size):
        memory[m] = (1 - gate[b,m]) * memory[m] + gate[b,m] * update[b]
\end{verbatim}

This approach had several problems:
\begin{itemize}
\item \textbf{Poor Performance:} Loops are much slower than vectorized operations
\item \textbf{No GPU Utilization:} Cannot take advantage of parallel processing
\item \textbf{Memory Inefficient:} Multiple small operations instead of large vectorized ones
\item \textbf{Not Scalable:} Performance degrades significantly with larger batch sizes
\end{itemize}

\subsection{Vectorized Solution}

The improved implementation uses batch-averaged vectorized operations:

\begin{equation}
\mathbf{M} \leftarrow (1 - \bar{\mathbf{g}} \mathbf{1}^T) \odot \mathbf{M} + (\bar{\mathbf{g}} \mathbf{1}^T) \odot (\mathbf{1} \bar{\mathbf{u}}^T) \quad \text{(broadcasting creates proper Hadamard dimensions)}
\end{equation}

where $\bar{\mathbf{g}} = \frac{1}{B}\sum_{b=1}^{B} \mathbf{g}[b, :]$ and $\bar{\mathbf{u}} = \frac{1}{B}\sum_{b=1}^{B} \mathbf{u}[b, :]$.

\subsection{Performance Results}

\begin{table}[h]
\centering
\begin{tabular}{|l|c|c|c|c|}
\hline
\textbf{Configuration} & \textbf{Simple} & \textbf{Attention} & \textbf{Sophisticated} \\
\hline
Small (1×10×64) & 0.136 ms & 0.617 ms & 0.603 ms \\
Medium (4×20×64) & 0.329 ms & 1.275 ms & 1.477 ms \\
Large (8×50×128) & 0.534 ms & 3.481 ms & 2.999 ms \\
XLarge (16×100×256) & 0.597 ms & 11.352 ms & 6.190 ms \\
\hline
\end{tabular}
\caption{Memory update performance (milliseconds per update)}
\end{table}

\subsection{Key Benefits}

\begin{itemize}
\item \textbf{Linear Scaling:} Performance scales much better with batch size
\item \textbf{GPU Acceleration:} Operations can run in parallel on GPU
\item \textbf{Memory Efficiency:} Single large operations instead of many small ones
\item \textbf{Code Quality:} Cleaner, more maintainable vectorized code
\end{itemize}

\appendix

\section{MultiHeadAttention Two-Array Return}

This appendix explains why the \texttt{nn.MultiheadAttention} function returns two arrays and how they are used in the memory update strategies.

\subsection{MultiHeadAttention Output Structure}

The \texttt{MultiHeadAttention} function returns two distinct outputs:

\begin{equation}
\mathbf{A}_{\text{mem}}, \mathbf{W}_{\text{attn},S} = \text{MultiHeadAttention}(\mathbf{x}, \mathbf{M}, \mathbf{M}) \in \mathbb{R}^{B \times S \times D} \times \mathbb{R}^{B \times S \times M}
\end{equation}

\subsection{Output 1: Attention Output ($\mathbf{A}_{\text{mem}}$)}

\textbf{Purpose:} Contains the memory-enhanced representation of the input.

\textbf{Dimension:} $\mathbb{R}^{B \times S \times D}$

\textbf{What it represents:} For each token in the input sequence $\mathbf{x}$, this contains a new representation that has integrated relevant information from memory $\mathbf{M}$.

\textbf{Usage:} This is the primary output that gets used as the update signal in the next step of the memory update process.

\subsection{Output 2: Attention Weights ($\mathbf{W}_{\text{attn},S}$)}

\textbf{Purpose:} Contains the attention scores showing how much each input token attends to each memory slot.

\textbf{Dimension:} $\mathbb{R}^{B \times S \times M}$

\textbf{What it represents:} For each batch item and sequence position, there's a vector of $M$ attention scores (one per memory slot).

\textbf{Usage:} These weights are used to determine which memory slots to update and by how much.

\subsection{Mathematical Breakdown}

The \texttt{MultiHeadAttention(x, M, M)} operation works as follows:

\begin{itemize}
\item \textbf{Query:} $\mathbf{x}$ (current state) - "What am I looking for?"
\item \textbf{Key:} $\mathbf{M}$ (memory buffer) - "What's available in memory?"
\item \textbf{Value:} $\mathbf{M}$ (memory buffer) - "What information should I retrieve?"
\end{itemize}

\textbf{Attention Computation:}
\begin{equation}
\text{Attention}(Q, K, V) = \text{softmax}\left(\frac{QK^T}{\sqrt{d_k}}\right)V
\end{equation}

Where:
\begin{itemize}
\item $Q = \mathbf{x}$ (queries)
\item $K = \mathbf{M}$ (keys)
\item $V = \mathbf{M}$ (values)
\end{itemize}

\subsection{Why Both Arrays Are Needed}

\subsubsection{For Memory Enhancement:}
\begin{itemize}
\item $\mathbf{A}_{\text{mem}}$ provides the enhanced input representation
\item This gets averaged across sequence length to create the update signal
\end{itemize}

\subsubsection{For Memory Update:}
\begin{itemize}
\item $\mathbf{W}_{\text{attn},S}$ provides the attention weights
\item These weights determine which memory slots were most relevant
\item Used to compute update gates: $\mathbf{g}_{\text{attn},\sigma} = \sigma(\text{mean}(\mathbf{W}_{\text{attn},S}))$
\end{itemize}

\subsection{Usage in Memory Update Strategies}

\subsubsection{Attention-Based Strategy:}
\begin{enumerate}
\item \textbf{Step 1:} Compute attention between input and memory
   \begin{equation}
   \mathbf{A}_{\text{mem}}, \mathbf{W}_{\text{attn},S} = \text{MultiHeadAttention}(\mathbf{x}, \mathbf{M}, \mathbf{M})
   \end{equation}

\item \textbf{Step 2:} Use $\mathbf{A}_{\text{mem}}$ to create update signal
   \begin{equation}
   \mathbf{s} = \mathbf{A}_{\text{mem}}.\text{mean}(\text{dim}=1) \quad \text{(Average across sequence)}
   \end{equation}

\item \textbf{Step 3:} Use $\mathbf{W}_{\text{attn},S}$ to create update gates
   \begin{equation}
   \mathbf{g}_{\text{attn},\sigma} = \sigma(\mathbf{W}_{\text{attn},S}.\text{mean}(\text{dim}=1)) \quad \text{(Attention-based gates)}
   \end{equation}

\item \textbf{Step 4:} Apply vectorized update
   \begin{equation}
   \mathbf{M} \leftarrow (1 - \bar{\mathbf{g}}_{\text{attn},\sigma} \mathbf{1}^T) \odot \mathbf{M} + (\bar{\mathbf{g}}_{\text{attn},\sigma} \mathbf{1}^T) \odot (\mathbf{1} \bar{\mathbf{s}}^T) \quad \text{(broadcasting creates proper Hadamard dimensions)}
   \end{equation}
   where $\bar{\mathbf{g}}_{\text{attn},\sigma} = \frac{1}{B}\sum_{b=1}^{B} \mathbf{g}_{\text{attn},\sigma}[b, :]$ and $\bar{\mathbf{s}} = \frac{1}{B}\sum_{b=1}^{B} \mathbf{s}[b, :]$
\end{enumerate}

\subsubsection{Sophisticated Strategy:}
The sophisticated strategy uses a similar pattern but with a single representative as the query:

\begin{enumerate}
\item \textbf{Step 1:} Compute attention between representative and memory
   \begin{equation}
   \mathbf{A}_{\text{soph}}, \mathbf{W}_{\text{soph},S} = \text{MultiHeadAttention}(\mathbf{r}_{\text{imp}}, \mathbf{M}, \mathbf{M})
   \end{equation}

\item \textbf{Step 2:} Use both outputs for sophisticated memory updates with decay and importance weighting
\end{enumerate}

\subsection{Key Insight}

The two outputs serve complementary roles:
\begin{itemize}
\item \textbf{$\mathbf{A}_{\text{mem}}$:} "What information was retrieved from memory?"
\item \textbf{$\mathbf{W}_{\text{attn},S}$:} "How much attention was paid to each memory slot?"
\end{itemize}

This dual output allows the attention-based strategy to both enhance the current state with memory information AND use the attention patterns to intelligently update the memory buffer itself.

\section{Architectural Design vs. Learned Behavior}

\subsection{The Gap Between Design Intent and Training Reality}

A critical consideration in the design of memory update mechanisms is the distinction between \textbf{what the architecture enables} and \textbf{what the model actually learns}. While the simple memory update strategy is designed to enable context-dependent learning through the formulation:

\begin{equation}
\mathbf{u} = \mathbf{W}_{\text{update}} \begin{bmatrix} \mathbf{r} \\ \bar{\mathbf{m}} \end{bmatrix}
\end{equation}

the actual behavior depends entirely on the learned weights $\mathbf{W}_{\text{update}} \in \mathbb{R}^{D \times 2D}$.

\subsection{The Hot Stove Analogy}

Consider the analogy of learning about fire safety through touching a hot stove:

\textbf{Baby's Experience:}
\begin{itemize}
\item Limited existing knowledge about heat, pain, and safety
\item Simple, direct association: ``hot stove → pain''
\item Memory update: Basic safety rule formation
\end{itemize}

\textbf{Teenager's Experience:}
\begin{itemize}
\item Rich existing knowledge about chemistry, physics, and safety protocols
\item Complex integration: ``hot stove → thermal energy transfer → cellular damage → pain response''
\item Memory update: Sophisticated understanding of heat transfer mechanisms
\end{itemize}

Both receive the same sensory input (touching the hot stove), but their memory integration strategies differ dramatically based on their existing knowledge base $\bar{\mathbf{m}}$.

\subsection{Theoretical vs. Practical Learning}

\textbf{Theoretical Capability:}
The architecture enables context-dependent updates through the concatenation $[\mathbf{r}, \bar{\mathbf{m}}]$, allowing the model to learn different integration strategies for different knowledge states.

\textbf{Practical Reality:}
The learned weights $\mathbf{W}_{\text{update}}$ are determined by:
\begin{itemize}
\item \textbf{Loss function design:} What behavior is rewarded during training
\item \textbf{Training data distribution:} What patterns the model observes
\item \textbf{Optimization success:} Whether the model finds optimal solutions
\end{itemize}

\subsection{Potential Learning Outcomes}

The model might learn:
\begin{itemize}
\item $\mathbf{W}_{\text{update}} \approx [\mathbf{W}_r, \mathbf{0}]$: Ignores memory context entirely
\item $\mathbf{W}_{\text{update}} \approx [\mathbf{0}, \mathbf{W}_m]$: Ignores new information entirely
\item $\mathbf{W}_{\text{update}} \approx [\mathbf{W}_r, \mathbf{W}_m]$: Uses both components (intended behavior)
\end{itemize}

\subsection{Empirical Validation Requirements}

To verify that context-dependent learning actually occurs, empirical analysis must examine:
\begin{enumerate}
\item \textbf{Weight analysis:} Whether $\mathbf{W}_{\text{update}}$ utilizes both input components
\item \textbf{Behavioral testing:} Whether updates vary with memory state
\item \textbf{Performance comparison:} Whether context-dependent updates improve task performance
\item \textbf{Ablation studies:} Whether removing memory context degrades performance
\end{enumerate}

\subsection{Design Implications}

This analysis highlights the importance of:
\begin{itemize}
\item \textbf{Careful loss function design} to encourage context-dependent behavior
\item \textbf{Diverse training data} that demonstrates context-dependent patterns
\item \textbf{Empirical validation} to verify that intended behaviors emerge
\item \textbf{Architectural constraints} that make desired behaviors more likely to be learned
\end{itemize}

The memory update mechanism provides the \textit{possibility} of sophisticated, context-dependent learning, but successful training is required to realize this potential.

\end{document}
